\subsection{Definition of polynomials}

Let I be a set. We recall (III, p.~452) that the free commutative algebra on I over A is denoted by \(\mathbf{A}[(\mathbf{X}_i)_{i\in I}]\) or \(\mathbf{A}[\mathbf{X}_i]_{i\in I}\) . The elements of this algebra are called polynomials with respect to the indeterminates \(\mathbf{X}_i\) (or in the indeterminates \(\mathbf{X}_i\) ) with coefficients in A. Let us recall that the indeterminate \(\mathbf{X}_i\) is the canonical image of \(i\) in the free commutative algebra on I over A; sometimes it is convenient to denote this image by another symbol such as \(\mathbf{X}_i'\) , \(\mathbf{Y}_i\) , \(\mathbf{T}_i\) , etc. This convention is often introduced by a phrase such as: « Let \(\mathbf{Y} = (\mathbf{Y}_i)_{i\in I}\) be a family of indeterminates »; in this case the algebra of polynomials in question is denoted by \(\mathbf{A}[\mathbf{Y}]\) . When \(\mathrm{I} = \{1,2,\dots,n\}\) , one writes \(\mathbf{A}[\mathbf{X}_1,\mathbf{X}_2,\dots,\mathbf{X}_n]\) in place of \(\mathbf{A}[(\mathbf{X}_i)_{i\in I}]\) .

For \(\nu \in \mathbf{N}^{(I)}\) we put

\[
\mathbf {X} ^ {\nu} = \prod_ {i \in I} \mathbf {X} _ {i} ^ {\nu_ {i}}.
\]

Then \((\mathbf{X}^{\nu})_{\nu \in \mathbf{N}^{(I)}}\) is a basis of the A-module \(\mathbf{A}[(\mathbf{X}_i)_{i\in I}]\) . The \(\mathbf{X}^{\nu}\) are called monomials in the indeterminates \(\mathbf{X}_i\) . For \(\nu = 0\) we obtain the unit element of \(\mathbf{A}[(\mathbf{X}_i)_{i\in I}]\) . Every polynomial \(u\in \mathbf{A}[(\mathbf{X}_i)_{i\in I}]\) can be written in exactly one way in the form

\[
u = \sum_ {\nu \in \mathbf {N} ^ {(I)}} \alpha_ {\nu} X ^ {\nu}
\]

where \(\alpha_{\nu}\in A\) and the \(\alpha_{\nu}\) are zero except for a finite number; the \(\alpha_{\nu}\) are called the coefficients of u; the \(\alpha_{\nu}X^{\nu}\) are called the terms of u (often the element \(\alpha_{\nu}X^{\nu}\) is called the term in \(X^{\nu}\) ), in particular the term of \(\alpha_{0}X^{0}\) , identified with \(\alpha_{0}\) , is called the constant term of u. When \(\alpha_{\nu}=0\) , we say by abuse of language that u contains no element in \(X^{\nu}\) ; in particular when \(\alpha_{0}=0\) , we say that u is a polynomial without constant term (III, p.~453). Any scalar multiple of 1 is called a constant polynomial.

Let B be a commutative ring and \(\rho: \mathbf{A} \to \mathbf{B}\) a ring homomorphism. We consider \(\mathrm{B}[(\mathbf{X}_i)_{i \in I}]\) as an A-algebra by means of \(\rho\) . Thus the mapping \(\sigma\) of \(\mathbf{A}[(\mathbf{X}_i)_{i \in I}]\) into \(\mathrm{B}[(\mathbf{X}_i)_{i \in I}]\) which transforms \(\sum \alpha_\nu \mathbf{X}^\nu\) into \(\sum \rho(\alpha_\nu) \mathbf{X}^\nu\) is a homomorphism of A-algebras; if \(u \in \mathbf{A}[(\mathbf{X}_i)_{i \in I}]\) , we sometimes denote by \(^p u\) the image of \(u\) by this homomorphism. The homomorphism of \(\mathbf{B} \otimes_{\mathbf{A}} \mathbf{A}[(\mathbf{X}_i)_{i \in I}]\) into \(\mathrm{B}[(\mathbf{X}_i)_{i \in I}]\) canonically defined by \(\sigma\) transforms, for every \(i \in I\) , \(1 \otimes \mathbf{X}_i\) into \(\mathbf{X}_i\) ; this is an isomorphism of B-algebras (III, p.~449).

Let M be a free A-module with basis \((e_{i})_{i\in I}\) . There exists precisely one unital homomorphism \(\varphi\) of the symmetric algebra \(\mathbf{S}(\mathbf{M})\) into the algebra \(\mathbf{A}[(\mathbf{X}_{i})_{i\in I}]\) such that \(\varphi(e_{i})=\mathbf{X}_{i}\) for each \(i\in I\) , and this homomorphism is an isomorphism (III, p.~506). This isomorphism is said to be canonical. It allows us to apply to polynomial algebras certain properties of symmetric algebras. For example, let \((\mathrm{I}_{\lambda})_{\lambda\in L}\) be a partition of I. Let \(\varphi_{\lambda}\) be the homomorphism of \(P_{\lambda}=\mathbf{A}[(\mathbf{X}_{i})_{i\in I_{\lambda}}]\) into \(P=\mathbf{A}[(\mathbf{X}_{i})_{i\in I}]\) which transforms \(X_{i}\) (qua element of \(P_{\lambda}\) ) into \(X_{i}\) (qua element of P). Then the \(\varphi_{\lambda}\) define a homomorphism of the algebra \(\otimes_{\lambda\in L}P_{\lambda}\) into the algebra P,

and this homomorphism is an isomorphism (III, p.~503, Prop. 9).

Let \(\mathbf{E}\) be an \(\mathbf{A}\) -module, and put \(\mathbf{E} \otimes_{\mathbf{A}} \mathbf{A}[(\mathbf{X}_i)_{i \in I}] = \mathbf{E}[(\mathbf{X}_i)_{i \in I}]\) . The elements of the \(\mathbf{A}\) -module \(\mathbf{E}[(\mathbf{X}_i)_{i \in I}]\) are called polynomials in the indeterminates \(\mathbf{X}_i\) with coefficients in \(\mathbf{E}\) . Such a polynomial can be written in just one way as \(\sum_{\nu \in \mathbf{N}^{(I)}} e_\nu \otimes \mathbf{X}^\nu\) , where \(e_\nu \in \mathbf{E}\) and the \(e_\nu\) are zero for all but a finite number of

suffixes; we frequently write \(e_{\nu}X^{\nu}\) instead of \(e_{\nu}\otimes X^{\nu}\) .

